\begin{abstract}
\noindent
Whether the six-sphere carries a complex structure is a classical open question. A manuscript circulated by Levent Alpöge claims one: a family of complex two-dimensional tori over the projective line, with local monodromies $T_1,T_2$ of orders $3$ and $4$, completed at its three special fibres, whose total space would be a compact complex threefold diffeomorphic to $S^6$ of algebraic dimension one. Such a threefold would contradict a theorem of Campana, Demailly and Peternell. The Zenodo record 10.5281/zenodo.22678442 verifies and reconstructs the manuscript. This paper sets out what the record establishes, with proofs where they are given here, and adds new results.

Two parts of the construction are determined completely here. The first is the step of the published proof that the record's correction note isolates. The paper audits the note and proves that the torsion responsible for the failure of the step is present exactly at the non-normal points of the fibre, and that at a reduced fibre with normal crossings the group which the step asserts to vanish is the space of sections of a line bundle on the normalization. The second is the integral period system. Its monodromy group is the kernel of an explicit character $\Phi$ of order $12$ on the stabiliser of an isotropic vector: a normal subgroup of index $12$ with cyclic quotient, a congruence subgroup of exact level $12$, arithmetic in its Zariski closure, and a non-split extension of $SL_2(\ZZ)$ by a Heisenberg group. The Levi half of $\Phi$ is, up to sign, the peripheral character of the record's higher-torus update. Its Heisenberg half is defined by the odd quadratic refinement of the mod-$2$ form, the odd spin structure of a torus, and on the standard Heisenberg subgroup it is the Heisenberg multiplier of the odd Jacobi theta function. The orbits of the monodromy group on the affine hyperplanes used by the Erdős--Straus workbench are classified at every level, and the exact value set of an octonionic cubic of the record's key advances is determined.

The paper also describes the bridges from the period system to paramodular forms, spin structures, degenerations of abelian surfaces and the author's other workbenches, with what is established about each. The existence of the threefold and the counterexample claim are not verified here.
\end{abstract}

\section*{Introduction}\phantomsection\label{sec:about}
\addcontentsline{toc}{section}{Introduction}

\subsection*{The problem and this record}

Campana, Demailly and Peternell proved that a compact complex threefold homeomorphic to $S^6$ would have algebraic dimension $0$ \cite[Corollary~2.3]{CDP20}. A complex structure on $S^6$ that fibres holomorphically over $\Pp^1$ would have algebraic dimension at least one, and would contradict that theorem. The manuscript \cite{AlpogeS6} claims such a fibration, built from the $(3,4,\infty)$ triangle group and a representation of it on $\ZZ^4$ (Section~\ref{sec:candidate}). The record, created by KokunoYumeto, with a verification by Codex and two review records by ChatGPT Pro, reconstructs the manuscript, repairs its local defects, and labels the construction and the counterexample as verified (Sections~\ref{sec:record} and~\ref{sec:ledger}). Its correction note isolates a step in the published proof that is not a formal consequence of the sequences displayed there; the manuscript had located the failure of that step at its own cusp fibre (Section~\ref{sec:correction}). Two later supplements build on the construction's data (Section~\ref{sec:later}). Philip Engel's exposition \cite{EngelS6} presents the construction with the stated aim of a self-contained proof.

This paper states what the record claims, proves what can be proved from the record's local, finite and algebraic statements, and adds the results below. It does not verify the global construction.

\subsection*{Main results}

The first result concerns the step. Let $S=\{g=0\}$ be a reduced hypersurface in a complex manifold $X$, with conormal sheaf $N^*$ and normalization $\eta:\tilde S\to S$; let $T_S$ be the torsion subsheaf of $\Omega^1_S$ and $\widetilde\Omega^1_S=\Omega^1_S/T_S$. The step asserts that $H^0(S,N^*\otimes A)=0$ and $H^0(S,\widetilde\Omega^1_S\otimes A)=0$ imply $H^0(S,\Omega^1_X|_S\otimes A)=0$.

\begin{introthm}[Propositions~\ref{prop:normal}, \ref{prop:kernel} and~\ref{prop:sharp}]
\textup{(i)} $T_S$ vanishes at $x$ if and only if $S$ is normal at $x$; so the step holds on normal fibres, and on a non-normal fibre it holds exactly when a connecting map $\delta_A$ is injective.
\textup{(ii)} If $S$ has normal crossings and $H^0(S,\widetilde\Omega^1_S\otimes A)=0$, then $H^0(S,\Omega^1_X|_S\otimes A)\cong H^0(\tilde S,\eta^*(N^*\otimes A))$. On a fibre, the conclusion of the step holds if and only if $H^0(\tilde S,\eta^*A)=0$.
\textup{(iii)} In the correction note's test family this group is one-dimensional for every value of the gluing parameter.
\end{introthm}

The normal case of (i) is the manuscript's Remark~10.4; the converse, (ii) and (iii) are new. So the step needs an extra hypothesis exactly at non-normal fibres, and (ii) states that hypothesis for the fibres with normal crossings that the construction uses.

The remaining results concern the integral period system. On $V=\ZZ^4$ with basis $(\gamma,u,w,\delta)$ the local monodromies $T_1,T_2$ fix $\gamma$ and preserve an alternating form $Q_0$ of type $(1,6)$. The stabiliser $G_{\ZZ}$ of $\gamma$ in $Sp(Q_0,\ZZ)$ is a semidirect product $SL_2(\ZZ)\ltimes\mathrm{Heis}(\ZZ)$, whose elements we write $L_Mh(p,r,s)$ (Lemma~\ref{lem:stab}). Let $\operatorname{ab}:SL_2(\ZZ)\to\ZZ/12$ be the abelianisation with $\operatorname{ab}\begin{psmallmatrix}1&1\\0&1\end{psmallmatrix}=1$, and put
\[
\Phi\bigl(L_Mh(p,r,s)\bigr)=\operatorname{ab}(M)+s-6(p^2+pr+r^2)\bmod12 .
\]

\begin{introthm}[Theorem~\ref{thm:mono}, Proposition~\ref{prop:twist}]
$\Phi$ is a surjective homomorphism, and the monodromy group $G=\langle T_1,T_2\rangle$ is its kernel. So $G$ is normal of index $12$ in $G_{\ZZ}$, with cyclic quotient generated by the central transvection $\delta\mapsto\delta+\gamma$. It maps onto $SL_2(\ZZ)$, it contains the principal congruence subgroup of level $12$ and none of level $4$ or $6$, and it is Zariski dense and arithmetic in the algebraic stabiliser of $\gamma$. The extension $1\to G\cap\mathrm{Heis}(\ZZ)\to G\to SL_2(\ZZ)\to1$ does not split.
\end{introthm}

\begin{introthm}[Proposition~\ref{prop:heis}, Theorem~\ref{thm:periph}]
\textup{(i)} The Heisenberg part $G\cap\mathrm{Heis}(\ZZ)$ is cut out by the odd quadratic refinement $p^2+pr+r^2$ of the mod-$2$ symplectic form, the only refinement invariant under $SL_2(\ZZ)$. On the standard Heisenberg subgroup it is the kernel of the character $(-1)^{p+r+pr+\kappa}$, the Heisenberg multiplier of the odd Jacobi theta function.
\textup{(ii)} The marked peripheral character $\chi$ of the record's higher-torus update equals $-\operatorname{ab}$ on the Levi parts, its values $(-1,4,-3)$ are those of $\operatorname{ab}$ on the Levi parts of $T_0,T_1,T_2$, and on $G$ it is determined by the Heisenberg coordinates: $\chi(M)\equiv s-6(p^2+pr+r^2)\pmod{12}$.
\end{introthm}

\begin{introthm}[Theorem~\ref{thm:orbitsS6}]
For every $D\ge1$, the orbits of $G$ on the affine hyperplane $H_D$ of the dual lattice modulo $D$ are the sets on which $\gcd(b,c,\gcd(D,6))$ is constant, where $b,c$ are the $\hat u$- and $\hat w$-coordinates, and they coincide with the orbits of $G_{\ZZ}$.
\end{introthm}

These results are new with respect to the record's files that were searched (\fn{00\_}--\fn{03\_}, \fn{14\_}--\fn{16\_} and \fn{26\_}): they use the period system, but contain no statement on the index, level, Zariski closure or arithmeticity of its monodromy group. The archive \fn{28\_} was not searched. The paper also proves the exact value set of the octonionic cubic of the key advance S6-5 (Proposition~\ref{prop:s65}) and a general form of the note's local computation of relative differentials (Proposition~\ref{prop:torsion}). By the latter, the record's correction to the manuscript's Remark~1.3 is right for relative one-forms, while relative two-forms do have torsion along the double curves of the cusp fibre.

\subsection*{What the results mean}

My assessment of the results on the step is that they settle its local content. The step is valid on normal fibres and needs an extra hypothesis on non-normal ones, and at the fibres with normal crossings that hypothesis is the vanishing of $H^0(\tilde S,\eta^*A)$. The global question, whether the hypotheses of \cite{CDP20} force this vanishing, is left open by the note. The record's answer at its own candidate is negative, and that answer depends on the construction, which is not verified here (Question~\ref{q:global}).

My assessment of the monodromy results is that they identify the peripheral character of the record's higher-torus update as the invariant that separates the monodromy group from its arithmetic hull, for the following reasons. The group is as large as it can be subject to one condition, $\Phi=0$. The Levi half of $\Phi$ is, up to sign, the update's character, and the Heisenberg half is the odd theta multiplier. So $G$ is exactly the locus where these two characters agree (Theorem~III(ii)). Because $G$ is arithmetic and not thin, the methods for congruence subgroups apply, and they give the orbit classification that the Erdős--Straus covers need at every level.

\subsection*{The bridges}

In the author's view the connections of this research to other areas are among its most valuable outcomes. The period system connects to several established theories, and Section~\ref{sec:overlaps} treats each with what is established and what is open:
\begin{itemize}
\item \emph{paramodular forms}: $Sp(Q_0,\ZZ)$ is the paramodular group of level $6$ of Gritsenko and Hulek, and $\Phi$ has the shape of one of their characters on it;
\item \emph{spin structures and theta characteristics}: the refinement that cuts out the Heisenberg part is the odd spin structure, which is also the Arf-one spin form of the update's Lie-framed torus;
\item \emph{degenerations}: the cusp monodromy has the Jordan type of Kulikov's type III on $\Lambda^2V$, and the manuscript's cusp fibre has the combinatorics of Mumford's degenerations of abelian surfaces;
\item \emph{thin groups}: the monodromy group is arithmetic, not thin.
\end{itemize}
The edges to the author's other workbenches, Yang--Mills, Erdős--Straus and Navier--Stokes, are listed with their status. The Erdős--Straus workbench uses exactly the dual action of Theorem~IV.

\subsection*{How this record was made}

The construction is that of the manuscript circulated by Levent Alpöge; the attribution file of the Yang--Mills workbench states that it was produced with Claude. The record is the work of KokunoYumeto, who designed the validate/repair/disprove programme and required proofs, calculations and counterexamples before any conclusion, with Codex for the verification and ChatGPT Pro for two independent reviews. Philip Engel's exposition is cited as \cite{EngelS6}.

This paper was prepared by Claude (Anthropic), model \texttt{claude-opus-5-5} (\emph{claude-ab} below), from an audit of the record on 26--27 September 2026. Independent Claude instances, each of which wrote none of the text, refereed the earlier version and the monodromy computations; they read for overstatement and for understatement and missed results, and several results were found in these passes, each credited at its statement. Every finding of another instance was verified by claude-ab before it was applied.

Every statement has one of four statuses:
\begin{itemize}
\item \emph{verified};
\item \emph{not verified here}, which implies nothing in either direction;
\item \emph{my assessment}, a view in the first person with its reasons, introduced by those words;
\item \emph{proved negative}, stated with its counterexample or derivation (Section~\ref{ssec:mononeg} collects those on the period system).
\end{itemize}
The status lines under the results make the first two precise. \textsf{Proved here} means that the proof is given in full in this paper; \textsf{Audited} means that the record's proof was read in full, every step checked, and independent checks run; \textsf{Reproduced} means that the finite content was recomputed and the general proof not re-derived. These are forms of \emph{verified}. \textsf{Located} means that the statement and its proof were found and are described, and \textsf{Stated} that the statement and the location of its proof are reported without the proof having been seen; both are forms of \emph{not verified here}. The record's own labels are \textsf{VERIFIED}, \textsf{OPEN} and \textsf{REFUTED}; they are reported as the record's, next to the status here.

The record's files were read as listed in Appendix~\ref{app:catalogue}: the correction note in full; the monograph's Part~I, ledger, Theorem~13.6 and §§13.11 and~14; the reconstructed manuscript's summary, Remarks~1.3 and~9.21, §2 (the period system) and §§10.1--10.3; the claim ledger; the guides and the key-advances reader. The rest was mapped from its tables of contents and summaries.

This is version~3. It keeps the theorems, proofs and citations of version~2 and corrects its framing, and it adds the integral monodromy of the period system (Section~\ref{sec:monodromy}) and its bridges. Version~2 and the audit notes cited as \note{NN\_} were withdrawn from the branch \texttt{claude/claude-ab-grind-20260925} on 27 September 2026 because of their framing; they remain, unchanged, in the branch's history at commit \texttt{a90d9e5b}, in \texttt{contrib/claude-ab/grind\_20260925/}. Appendix~\ref{app:changes} lists every change.

\subsection*{Organization}

Section~\ref{sec:problem} states the question and what the record claims, Section~\ref{sec:record} maps the record, Section~\ref{sec:candidate} outlines the construction, and Section~\ref{sec:ledger} reports the record's ledger. Section~\ref{sec:correction} audits the correction note and proves Theorem~I. Section~\ref{sec:later} treats the later supplements and the arithmetic of S6-5. Section~\ref{sec:monodromy} determines the monodromy group and proves Theorems~II--IV. Section~\ref{sec:overlaps} describes the bridges, and Section~\ref{sec:open} the open questions and what was not checked. The appendices give the catalogue of statements, the verification and the changes of version~3.
